\title{GQA: Training Generalized Multi-Query Transformer Models from Multi-Head Checkpoints}

\begin{abstract}
Multi-query attention (\textsc{MQA}), which only uses a single key-value head, drastically speeds up decoder inference. However, \textsc{MQA} can lead to quality degradation, and moreover it may not be desirable to train a separate model just for faster inference. We (1) propose a recipe for uptraining existing multi-head language model checkpoints into models with \textsc{MQA} using 5\% of original pre-training compute, and (2) introduce grouped-query attention (\textsc{GQA}), a generalization of multi-query attention which uses an intermediate (more than one, less than number of query heads) number of key-value heads. We show that uptrained \textsc{GQA} achieves quality close to multi-head attention with comparable speed to \textsc{MQA}.
\end{abstract}

\section{Introduction}
\label{section:intro}

Autoregressive decoder inference is a severe bottleneck for Transformer models due to the memory bandwidth overhead from loading decoder weights and all attention keys and values at every decoding step~\citep{shazeer2019mq, palminference, dejong2022fido}. The memory bandwidth from loading keys and values can be sharply reduced through \emph{multi-query attention}~\citep{shazeer2019mq}, which uses multiple query heads but single key and value heads.

However, multi-query attention (\textsc{MQA}) can lead to quality degradation and training instability, and it may not be feasible to train separate models optimized for quality and inference. Moreover, while some language models already use multi-query attention, such as PaLM~\citep{palm}, many do not, including publicly available language models such as T5~\citep{t5} and LLaMA~\citep{llama}. 

This work contains two contributions for faster inference with large language models. First, we show that language model checkpoints with multi-head attention (\textsc{MHA}) can be \emph{uptrained}~\citep{sparsemoe} to use \textsc{MQA} with a small fraction of original training compute. This presents a cost-effective method to obtain fast multi-query as well as high-quality \textsc{MHA} checkpoints.

Second, we propose grouped-query attention (\textsc{GQA}), an interpolation between multi-head and multi-query attention with single key and value heads \emph{per subgroup of query heads}. We show that uptrained \textsc{GQA} achieves quality close to multi-head attention while being almost as fast as multi-query attention.

\section{Method}
\label{section:method}

\subsection{Uptraining}

Generating a multi-query model from a multi-head model takes place in two steps: first, converting the checkpoint, and second, additional pre-training to allow the model to adapt to its new structure. Figure \ref{fig:recycling} shows the process for converting a multi-head checkpoint into a multi-query checkpoint. The projection matrices for key and value heads are mean pooled into single projection matrices, which we find works better than selecting a single key and value head or randomly initializing new key and value heads from scratch.

The converted checkpoint is then pre-trained for a small proportion $\alpha$ of its original training steps on the same pre-training recipe.

\subsection{Grouped-query attention}

Grouped-query attention divides query heads into $G$ \emph{groups}, each of which shares a single key head and value head. \textsc{GQA}-\textsc{g} refers to grouped-query with $G$ groups. \textsc{GQA}-$1$, with a single group and therefore single key and value head, is equivalent to \textsc{MQA}, while \textsc{GQA}-\textsc{h}, with groups equal to number of heads, is equivalent to \textsc{MHA}. Figure \ref{fig:gmq_architecture} shows a comparison of grouped-query attention and multi-head/multi-query attention. When converting a multi-head checkpoint to a \textsc{GQA} checkpoint, we construct each group key and value head by mean-pooling all the original heads within that group.

An intermediate number of groups leads to an interpolated model that is higher quality than \textsc{MQA} but faster than \textsc{MHA}, and, as we will show, represents a favorable trade-off. Going from \textsc{MHA} to \textsc{MQA} reduces $H$ key and value heads to a single key and value head, reducing the size of the key-value cache and therefore amount of data that needs to be loaded by a factor of $H$. However, larger models generally scale the number of heads, such that multi-query attention represents a more aggressive cut in both memory bandwidth and capacity. \textsc{GQA} lets us keep the same proportional decrease in bandwidth and capacity as model size increases.

Moreover, larger models suffer relatively less from memory bandwidth overhead from attention, as the KV-cache scales with model dimension while model FLOPs and parameters scale with the square of model dimension. Finally, standard sharding for large models replicates the single key and value head by the number of model partitions~\citep{palminference}; \textsc{GQA} removes the waste from such partitioning. Therefore, we expect \textsc{GQA} to present a particularly good trade-off for larger models. 

We note that \textsc{GQA} is not applied to the encoder self-attention layers; encoder representations are computed in parallel, and memory bandwidth is therefore generally not the primary bottleneck.

\section{Experiments}
\label{section:experiments}

\subsection{Experimental setup}

\paragraph{Configurations}
All models are based on the T5.1.1 architecture~\citep{t5}, implemented with JAX~\citep{jax}, Flax~\citep{flax}, and Flaxformer\footnote{\url{https://github.com/google/flaxformer}}. For our main experiments we consider T5 Large and XXL with multi-head attention, as well as uptrained versions of T5 XXL with multi-query and grouped-query attention. We use the Adafactor optimizer with the same hyperparameters and learning rate schedule as T5~\citep{t5}. We apply \textsc{MQA} and \textsc{GQA} to decoder self-attention and cross-attention, but not encoder self-attention.

\paragraph{Uptraining}

Uptrained models are initialized from public T5.1.1 checkpoints. The key and value heads are mean-pooled to the appropriate \textsc{MQA} or \textsc{GQA} structure, and then pre-trained for a further 
$\alpha$ proportion of original pre-training steps with the original pre-training setup and dataset from~\citep{t5}. For $\alpha=0.05$, training took approximately 600 TPUv3 chip-days.

\paragraph{Data}
We evaluate on summarization datasets CNN/Daily Mail~\citep{cnn}, arXiv and PubMed~\citep{arxiv}, MediaSum~\citep{mediasum}, and Multi-News~\citep{multinews}; translation dataset WMT 2014  English-to-German; and question answering dataset TriviaQA~\citep{triviaqa}. We do not evaluate on popular classification benchmarks such as GLUE~\citep{glue} as autoregressive inference is less applicable for those tasks. 

\paragraph{Fine-tuning}

For fine-tuning, we use a constant learning rate of 0.001, batch size 128, and dropout rate 0.1 for all tasks. CNN/Daily Mail and WMT use input length of 512 and output length 256. Other summarization datasets use input length 2048 and output length 512. Finally, TriviaQA uses input length 2048 and output length 32. We train until convergence and select the checkpoint with the highest dev performance. We use greedy decoding for inference.

\paragraph{Timing}
We report time per sample per TPUv4 chip, as measured by xprof~\citep{xprof}. For timing experiments we use 8 TPUs with the largest batch size that fits up to 32 per TPU, and parallelization optimized separately for each model.

\subsection{Main results}

Figure \ref{fig:perf_vs_time} shows average performance over all datasets as a function of average inference time for \textsc{MHA} T5-Large and T5-XXL, and uptrained \textsc{MQA} and \textsc{GQA}-$8$ XXL models with uptraining proportion $\alpha = 0.05$. We see that a larger uptrained \textsc{MQA} model provides a favorable trade-off relative to \textsc{MHA} models, with higher quality and faster inference than \textsc{MHA}-Large. Moreover, \textsc{GQA} achieves significant additional quality gains, achieving performance close to \textsc{MHA}-XXL with speed close to \textsc{MQA}. Table \ref{table:headline_results} contains full results for all datasets.

\subsection{Ablations}

This section presents experiments to investigate the effect of different modeling choices. We evaluate performance on a representive subsample of tasks: CNN/Daily Mail, (short-form summarization), MultiNews (long-form summarization), and TriviaQA (question-answering).

\paragraph{Checkpoint conversion} Figure \ref{fig:conversion_methods} compares the performance of different methods for checkpoint conversion. Mean pooling appears to work best, followed by selecting a single head and then random initialization. Intuitively, results are ordered by the degree to which information is preserved from the pre-trained model.

\paragraph{Uptraining steps} Figure \ref{fig:uptraining_steps} shows how performance varies with uptraining proportion for T5 XXL with \textsc{MQA} and \textsc{GQA}. First, we note that \textsc{GQA} already achieves reasonable performance after conversion while \textsc{MQA} requires uptraining to be useful. Both \textsc{MQA} and \textsc{GQA} gain from 5\% uptraining with diminishing returns from 10\%.

\paragraph{Number of groups}

Figure \ref{fig:time_groups} demonstrates the effect of the number of \textsc{GQA} groups on inference speed. For larger models the memory bandwidth overhead from the KV cache is less constraining~\citep{shazeer2019mq}, while the reduction in key-value size is sharper due to the increased number of heads. As a result, increasing the number of groups from \textsc{MQA} only results in modest slowdowns initially, with increasing cost as we move closer to \textsc{MHA}. We selected 8 groups as a favorable middle ground.

\section{Related Work}

This work is focused on achieving a better trade-off between decoder quality and inference time through reducing the memory bandwidth overhead~\citep{roofline} from loading keys and values. \citet{shazeer2019mq} first proposed reducing this overhead through multi-query attention. Follow-up work showed that multi-query attention is especially helpful for long inputs~\citep{palminference, dejong2022fido}. \citet{gqamrabe} independently developed GQA with public implementation. Other works have explored grouping attention heads for computational efficiency~\citep{grouptrans, gmha, pillars} without focusing specifically on key-value heads, which determine memory bandwidth overhead.

A number of other methods have been proposed to reduce memory bandwidth overhead from keys and values, as well as parameters. Flash attention~\citep{flashattention} structures the attention computation to avoid materializing the quadratic attention scores, reducing memory and speeding up training. Quantization~\citep{int8, gptq} reduces the size of weights and activations, including keys and values, by lowering precision. Model distillation~\citep{distill, distillsurvey} instead reduces model size at a given precision, using data generated from the larger model to finetune the smaller model. Layer-sparse cross-attention \citep{dejong2022fido} eliminates most cross-attention layers which make up the primary expense for longer inputs. Speculative sampling~\citep{specchen, specleviathan} ameliorates the memory bandwidth bottleneck by proposing multiple tokens with a smaller model which are then scored in parallel by a larger model. 

Finally, the uptraining procedure we propose is inspired by \citet{sparsemoe}, which uptrains standard T5 checkpoints into sparsely activated Mixture-of-Experts models.

\label{section:related}

\section{Conclusion}
\label{section:conclusion}

Language models are expensive for inference primarily due to the memory bandwidth overhead from loading keys and values. Multi-query attention reduces this overhead at the cost of decreased model capacity and quality. We propose to convert multi-head attention models to multi-query models with a small fraction of original pre-training compute. Moreover, we introduce grouped-query attention, an interpolation of multi-query and multi-head attention that achieves quality close to multi-head at comparable speed to multi-query attention. 

\section*{Limitations}

This paper focuses on ameliorating the memory bandwidth overhead from loading keys and values. This overhead is most important when generating longer sequences, for which quality is inherently difficult to evaluate. For summarization we employ Rouge score, which we know is a flawed evaluation that does not tell the whole story; for that reason, it is difficult to be certain our trade-offs are correct. Due to limited computation, we also do not compare our XXL \textsc{GQA} model to a comparitive model trained from scratch, so we do not know the relative performance of uptraining vs training from scratch. Finally, we evaluate the impact of uptraining and \textsc{GQA} only on encoder-decoder models. Recently, decoder-only models are extremely popular, and since these models do not have separate self-attention and cross-attention, we expect \textsc{GQA} to have a stronger advantage over \textsc{MQA}.

\section*{Acknowlegements}
We thank Santiago Onta\~{n}\'{o}n, Afroz Mohiuddin, William Cohen and others at Google Research for insightful advice and discussion. 

\end{document}
